\section{An actual smoothed map on the pole-free common-sum domain}
\label{sec:pole-free-smoothed-transport}

The raw iteration obstruction in MA-MO-008 does not preclude a
bounded once-defined smoothed map. We now construct such a map on
the exact pole-free source domain, with its inherited two-sided
Hilbert norm. The result is a trace-class map between two specified
Hilbert spaces. It does not assert an endomorphism of the smaller
domain or assign it an endomorphism trace.

Keep the rational source place set \(\{\infty\}\), the real source
cutoff \(\phi\), its additive version \(v\), and
\(L=1-M_v-\mathfrak F M_v\mathfrak F\) from MA-MO-008.
Thus \(\phi=1\) sufficiently far left and \(\phi=0\)
sufficiently far right. For \(0<a<1/2\) and \(s\ge0\), define
\[
 \begin{aligned}
 p_s(x)&=(1+x^2)^{s/2},\\
 H_{b,s}&=L^2(\mathbb R,e^{2bx}p_s(x)\,dx)\quad(b=\pm a),\\
 H_{\Sigma,a,s}
   &=L^2\bigl(\mathbb R,(e^{2ax}+e^{-2ax})p_s(x)\,dx\bigr),\\
 \mathscr Y_\phi&=\Theta\mathscr U_{\rm pf},\qquad
 K_{a,s,\phi}=\overline{\mathscr Y_\phi}^{\,H_{\Sigma,a,s}},\qquad
 J_{a,s}=\overline{\Theta\mathscr U_0}^{\,H_{\Sigma,a,s}}.
 \end{aligned}
 \tag{MA-MO.120}
\]
The source spaces \(\mathscr U_0,\mathscr U_{\rm pf}\) are
exactly (MA-MO.89), including both cutoff moment conditions.
The theta map is injective on \(\mathscr U_0\), by the
actual theta--M\"obius identity MA-TR-101. Consequently the
following action on \(\mathscr Y_\phi\) is well defined:
\[
       B_{\phi,g}^{\rm raw}(\Theta f)=C_g\Theta Lf,
                   \qquad f\in\mathscr U_{\rm pf}.
 \tag{MA-MO.121}
\]
It is the first component of the original
\(C_g(P_+-P_-)i_+f\); its second component is the same.

There are two analytic regimes in the theorem:
\begin{enumerate}
\item \emph{Fixed horizontal gap.} For some \(\delta>0\),
\[
 \left|\tfrac12+a-\operatorname{Re}\rho\right|\ge\delta
                  \quad\hbox{for every nontrivial zeta zero }\rho.
 \tag{MA-MO.122}
\]
Here \(g\in C_c^\infty(\mathbb R)\) is arbitrary and
\(\phi'\in C_c^\infty(\mathbb R)\). Under the stipulated
input A for rational zeta, this regime holds at every fixed
\(0<a<1/2\), with \(\delta=a\).
\item \emph{Unconditional generic width.} The number
\(\tfrac12+a\) belongs to the full-measure set
\(\mathscr G\) of MA-RS-001. Both \(g\) and \(\phi'\)
are compactly supported of some finite Gevrey order \(q>1\).
No zero-location hypothesis is used in this regime.
\end{enumerate}

\paragraph{MA-MO-010 (actual pole-free smoothing in every polynomial weight).}
In either stated regime, (MA-MO.121) extends uniquely to a
trace-class map
\[
             B_{\phi,g}:K_{a,s,\phi}\longrightarrow J_{a,s}
                    \subseteq H_{\Sigma,a,s}.
 \tag{MA-MO.123}
\]
Thus the original smoothed matrix, restricted to its genuine
pole-free smooth diagonal, has a bounded extension into the
same ambient narrow Hilbert product, for every \(s\ge0\).
Its two components are obtained from compatible one-sided
extensions, not by replacing the sum norm by a one-sided norm.

For fixed \(a,s,\phi\), the map \(g\mapsto B_{\phi,g}\)
is continuous into trace class on each fixed compact smooth
test-support space in the gap regime. In the generic regime
it is uniformly bounded for families with fixed support and
common finite-order Gevrey bounds. No invariance of
\(K_{a,s,\phi}\) under this map is asserted.

\begin{proof}
\emph{The estimates on both lines.}
Put
\[
 A_b(t)=\zeta(\tfrac12+b+it),\qquad j_b(t)=A_b(t)^{-1},
 \qquad G_b(t)=\mathcal Mg(b+it),\qquad b=\pm a.
 \tag{MA-MO.124}
\]
The nontrivial zero set is invariant under
\(\rho\mapsto1-\overline\rho\). Consequently (MA-MO.122)
gives the same fixed horizontal gap for \(\tfrac12-a\).
Likewise (MA-RS.2) for \(\tfrac12+a\) gives precisely the
same bound for \(\tfrac12-a\), since that reflection keeps
the height \(\operatorname{Im}\rho\) unchanged.

The minimum-modulus proofs of MA-RS-001 and MA-RS-004 apply
on both lines. To check their scope, their disk is centered
at \(2+it\), has radius between 12 and 13, and contains both
fixed real parts in its inner quarter. At large height it
contains no zeta pole or trivial zero. The same local zero
count, polynomial upper bound, and Harnack estimate apply.
The reflected distance bound supplies the same lower bound
for each finite Blaschke factor. At bounded height the line
has no zero, and its real part lies strictly between 0 and 1.
Thus, for every fixed derivative order, \(j_b\) has
polynomial growth in the gap regime, and is bounded by
\(\exp(C\log^2(2+|t|))\) in the generic regime. All
derivatives of \(A_b\) have polynomial growth in both cases.
This explicitly verifies the left-line estimate; it is not
deduced merely from a zero-free line without a rate.

\emph{The two one-sided smoothing kernels.}
Let \(V_b=\mathcal F_+M_{e^{bx}}\), as in MA-RS.8, and
let \(Z_b,J_b\) denote the maximal multiplier operators
with symbols \(A_b,j_b\). Initially on compact-frequency
tests define the sequential difference
\(D_b=M_\phi-Z_bM_\phi J_b\). The RS distributional
calculation gives for \(C_gD_b\) the kernel
\[
 \begin{aligned}
 k_{b,g}(t,u)
    &=\frac{i}{2\pi}G_b(t)A_b(t)
             \frac{j_b(t)-j_b(u)}{t-u}\,P(t-u),\\
 P(w)&=\int e^{iwx}\phi'(x)\,dx.
 \end{aligned}
 \tag{MA-MO.125}
\]
It is a Schwartz kernel. In the gap regime this follows
from the polynomial derivative estimates, the integral
formula for the divided difference, and the Schwartz
functions \(G_b(t)\) and \(P(t-u)\), exactly as in
MA-RS.37. In the generic regime their Gevrey Fourier decay
beats \(\exp(C\log^2(2+|t|+|u|))\), as explicitly proved
in MA-RS.22--.23. These arguments apply to every mixed
derivative. Denote the resulting operator by \(K_b(g)\).

Here is why the same kernel gives trace class in every
polynomial weight, not only the exponential weight.
Let \(r=s/2\) and
\(\Lambda^r=(1-\partial_t^2)^{r/2}\). Plancherel gives
\[
   \|u\|_{H_{b,s}}
      =\|\Lambda^r V_bu\|_2.
 \tag{MA-MO.126}
\]
Thus \(V_b\) identifies this space with the Sobolev space
\(H^r(\mathbb R)\). Conjugating a Schwartz-kernel operator
by \(\Lambda^r\) gives the kernel
\(\Lambda_t^r\Lambda_u^{-r}k_{b,g}(t,u)\). It remains
Schwartz: after Fourier transformation in the two variables,
the multipliers are \((1+\xi^2)^{r/2}\) and
\((1+\eta^2)^{-r/2}\), whose derivatives have polynomial
growth. The Hermite expansion argument of MA-RS-002 now
proves trace class on ordinary \(L^2\), hence on
\(H_{b,s}\) before conjugation. The estimates use finitely
many Schwartz seminorms and give the claimed test-family
bounds. This also proves the corresponding statement for
any other Schwartz kernel used below.

\emph{Agreement on the actual theta domain.}
The additive source need not be flat at zero. If
\(f\in\mathscr U_0\), define
\[
 F(x)=2e^{x/2}f(e^x),\qquad
 \widetilde F(x)=2e^{x/2}\mathfrak Ff(e^x),\qquad
                         y=\Theta f.
 \tag{MA-MO.127}
\]
Both weighted functions \(e^{bx}F\) and
\(e^{bx}\widetilde F\) are Schwartz for \(b=\pm a\).
Indeed evenness and the vanishing values of \(f\) and
\(\mathfrak Ff\) at zero give \(O(t^2)\) there,
with the corresponding bounds for all logarithmic
derivatives. The radial functions therefore decrease like
\(e^{5x/2}\) at the left, while their right tails decrease
faster than every exponential. This suffices for all the
weights presently used; it does not assert \(F\in\mathcal E\).

The actual theta sums are
\(Z^{\rm geom}F=y\) and
\(Z^{\rm geom}\widetilde F=Ry\), by Poisson summation.
Their moments at \(1/2\) are zero. The contour argument
of MA-MO-004 applies here on the strip between a fixed
right line \(b'>1/2\) and either \(b=\pm a\): the
Mellin transforms of these radial functions are holomorphic
for \(\Re z>-5/2\), with rapid vertical decay on each
smaller closed strip. No nonzero pole residue remains.
Consequently
\[
        V_by=A_bV_bF,\qquad
        V_bRy=A_bV_b\widetilde F,
        \qquad J_by=F,\quad J_bRy=\widetilde F.
 \tag{MA-MO.128}
\]
These are identities in the actual maximal inverse domains.
The same reasoning applies to \(\phi F\) and
\(\phi\widetilde F\). For \(f\in\mathscr U_{\rm pf}\)
their two moments vanish exactly, so
\[
   Z_b(\phi F)=\Theta(vf),\qquad
   Z_b(\phi\widetilde F)=\Theta(v\mathfrak Ff).
 \tag{MA-MO.129}
\]
Every output in this display is in \(\mathcal E\).

We also verify that the bounded kernel extension agrees
with these sequential outputs. Put \(u=V_by\) and
\(q=j_bu=V_bF\), both Schwartz. Let \(\chi_N\) be
standard smooth frequency cutoffs tending to one, and
put \(u_N=\chi_Nu\). Then
\(j_bu_N=\chi_Nq\to q\) in Schwartz space, and
\(u_N\to u\) there. Multiplication by \(\phi\) in
physical space and by \(A_b\) in frequency space preserve
Schwartz space. Hence the sequential outputs converge
in Schwartz topology, thus in every Sobolev norm.
The kernel extensions are continuous in those norms and
agree on \(u_N\), so they agree on \(u\). This argument
also works for \(Ry\). In particular it does not require
\(j_b\) to preserve every Schwartz function in the
generic regime; it uses the exact cancellation
\(j_bu_N=\chi_Nq\).

\emph{The original two-component action and the compact residual.}
For \(f\in\mathscr U_{\rm pf}\), both \(vf\) and
\(v\mathfrak Ff\) have value and integral zero.
Poisson summation therefore gives
\[
 \begin{aligned}
 \Theta Lf
 &=y-\Theta(vf)-R\Theta(v\mathfrak Ff)\\
 &=\psi y+D_by+RD_{-b}Ry,\qquad
                   \psi(x)=1-\phi(x)-\phi(-x).
 \end{aligned}
 \tag{MA-MO.130}
\]
The function \(\psi\) is compactly supported and smooth.
It is identically zero for a symmetric source cutoff.
Reflection is unitary from \(H_{b,s}\) to \(H_{-b,s}\),
because \(p_s\) is even, and
\(C_gR=RC_{g^\vee}\), where \(g^\vee(x)=g(-x)\).
It follows that the actual output in (MA-MO.121) agrees,
on the entire stated smooth core, with the trace-class
operator on \(H_{b,s}\)
\[
 T_b(g)=C_gM_\psi+K_b(g)+RK_{-b}(g^\vee)R.
 \tag{MA-MO.131}
\]
For completeness, the weighted Fourier kernel of the first
term is
\(G_b(t)(2\pi)^{-1}\int e^{i(t-u)x}\psi(x)\,dx\).
It is Schwartz, since both factors have rapid decay in
the independent variables \(t,t-u\). Thus the polynomial
weight proof above applies to it too. No nonsymmetric
cutoff term has been omitted.

\emph{Compatibility in the inherited sum norm.}
Set \(\mathbb H=H_{a,s}\oplus H_{-a,s}\) and define
the isometric embedding
\[
 \iota:H_{\Sigma,a,s}\longrightarrow\mathbb H,
             \qquad \iota y=(y,y),\qquad
 \mathbb T_g=\operatorname{diag}(T_a(g),T_{-a}(g)).
 \tag{MA-MO.132}
\]
The range of \(\iota\) is closed. The operator
\(\mathbb T_g\) is trace class, and by the actual-domain
identity just proved,
\[
       \mathbb T_g\iota y=\iota B_{\phi,g}^{\rm raw}y
                       \qquad(y\in\mathscr Y_\phi).
 \tag{MA-MO.133}
\]
Thus, if such \(y_n\) converge to \(y\in K_{a,s,\phi}\),
their output pairs converge in \(\mathbb H\) and their
limit is still in the closed diagonal range. Define
\[
 B_{\phi,g}=\left.\iota^*\mathbb T_g\iota\right|_{K_{a,s,\phi}}.
 \tag{MA-MO.134}
\]
On that domain this is precisely the inverse diagonal
identification of the output pair, by (MA-MO.133) and
closedness. It does not average two incompatible outputs.
Composition with bounded maps proves it is trace class,
with
\[
       \|B_{\phi,g}\|_1
           \le\|T_a(g)\|_{1,H_{a,s}}
                    +\|T_{-a}(g)\|_{1,H_{-a,s}}.
 \tag{MA-MO.135}
\]
Uniqueness follows from density of \(\mathscr Y_\phi\)
in its defined closure.

The range is in \(J_{a,s}\), not just in the ambient
space. On the smooth core \(Lf\in\mathscr U_0\), and
\[
 C_g\Theta Lf=\Theta h_g,\qquad
 h_g(t)=\int g(z)e^{-z/2}(Lf)(e^{-z}t)\,dz.
 \tag{MA-MO.136}
\]
Compactness of the dilation parameter gives an even
Schwartz function; its value and integral are zero.
Taking closure proves the range assertion. Restriction
of the codomain to this closed subspace preserves the
singular values. All trace-norm continuity and uniformity
assertions follow from the earlier kernel bounds and
the bounded compositions in (MA-MO.134).
\end{proof}

\paragraph{What has and has not been transported.}
The action on the smooth core is the original smoothed
common-sum matrix, with its two additional polar conditions
verified before any Hilbert norm is taken. In particular
this is an actual source identification, beyond a formal
comparison of scalar traces or of continued multipliers.
Both components use the same inherited sum norm, including
every polynomial exponent \(s\ge0\).

The closed domain \(K_{a,s,\phi}\) has been defined as a
closure, not identified with \(J_{a,s}\) or with an
arbitrarily chosen finite-codimensional Hilbert subspace.
No claim that \(B_{\phi,g}(K_{a,s,\phi})\subseteq
K_{a,s,\phi}\) is made. Trace class between two Hilbert
spaces does not itself supply an endomorphism trace,
a determinant, or an iterated original source action.
MA-MO-008's raw second-iteration obstruction remains valid.
The result also does not identify an orthogonal global
Paley--Wiener projection with a local Fourier cutoff.
No continuous lift from the Hilbert completion into the
original all-weight smooth space is inferred merely from
these fixed-weight Schwartz kernels.

The generic-width regime uses the proved RS distance
estimate and Gevrey tails, and covers a full-measure
set of widths, not every width. The arbitrary-smooth
regime uses the explicit fixed horizontal gap; invoking
A is one sufficient source of that gap. The arithmetic
polynomial cost in MA-PR-010/011 is not proved by either
regime. The inherited cote~1 locators, scans~079--080 and
111--112 via BR-TR-01/02, apply to actual nuclear maps
with their stated domains; finite-length spectral data
do not turn the map (MA-MO.123) into an endomorphism.
